\section{Background}
\label{sec:background}
\subsection{Bài toán N-Queens và điều kiện không tấn công}
Đánh số hàng, cột từ $0$ đến $N-1$. Hai quân hậu tại $(i,j)$ và $(i',j')$ tấn công nhau khi
\begin{equation}
i=i'\ \lor\ j=j'\ \lor\ |i-i'|=|j-j'|.
\label{eq:attack}
\end{equation}
Điều kiện đường chéo có thể được tách thành $i-j=i'-j'$ hoặc $i+j=i'+j'$. Vì cần đặt đúng $N$ quân hậu và mỗi hàng có tối đa một hậu, mỗi hàng phải có đúng một hậu. Tương tự đối với cột. Đường chéo chỉ cần tối đa một hậu, không cần chứa hậu.

Một nghiệm có thể được biểu diễn bằng vector $q=(q_0,\ldots,q_{N-1})$, trong đó $q_i$ là cột chứa hậu ở hàng $i$. Việc kiểm tra nghiệm dựa trên:
\begin{equation}
\begin{gathered}
0\le q_i<N\quad(0\le i<N),\\
|\{q_i\}|=|\{i+q_i\}|=|\{i-q_i\}|=N.
\end{gathered}
\label{eq:check}
\end{equation}
Ba tập có lực lượng $N$ nghĩa là không trùng cột và không trùng đường chéo ở cả hai hướng. Vector $(2,0,3,1)$ là ví dụ hợp lệ khi $N=4$.

\subsection{Logic mệnh đề, literal, clause và CNF}
Một biến Boolean nhận True hoặc False. Literal là một biến $x$ hoặc phủ định $\neg x$. Một clause là phép OR của các literal; công thức dạng chuẩn hội (CNF) là phép AND của các clause. Ví dụ
\begin{equation}
(x_1\lor x_2\lor x_3)\land(\neg x_1\lor\neg x_2)
\end{equation}
có ba biến, hai clause và năm lần xuất hiện literal. Số literal được đếm theo số lần xuất hiện, không phải số biến phân biệt.

Một công thức là SAT nếu tồn tại phép gán làm tất cả clause đúng; UNSAT nếu không tồn tại. Ví dụ trên cho phép một số phép gán thỏa, trong khi $x\land\neg x$ là UNSAT. Hai quy tắc chuyển đổi sử dụng nhiều trong mã hóa là
\begin{align}
A\Rightarrow B&\equiv\neg A\lor B,\\
A\Rightarrow(B\land C)&\equiv
(\neg A\lor B)\land(\neg A\lor C).
\end{align}
Trong giao diện PySAT, ID dương biểu diễn biến và ID âm biểu diễn phủ định. List \texttt{[-1,17]} là clause $\neg x_1\lor y_{17}$. Các list clause được hội lại thành CNF. Trong file DIMACS, cùng clause được ghi \texttt{-1 17 0}; số 0 là dấu kết thúc, không phải biến hoặc literal.

\subsection{ALO, AMO và Exactly-One}
Với $X=(x_1,\ldots,x_k)$:
\begin{align}
\ALO(X)&=\bigvee_{t=1}^{k}x_t,\\
\AMO(X)&:\ \sum_{t=1}^{k}[x_t]\le1,\\
\EO(X)&=\ALO(X)\land\AMO(X).
\label{eq:eo}
\end{align}
Ở đây $[x_t]=1$ khi $x_t$ đúng, ngược lại bằng 0. ALO loại phép gán tất cả sai; AMO loại các phép gán có ít nhất hai biến đúng. AMO riêng lẻ phải cho phép không có biến nào đúng.

\begin{table}[pos=H]\centering\small
\caption{Ý nghĩa các ràng buộc lực lượng trên một nhóm biến.}
\begin{tabular}{@{}lccc@{}}\toprule
Số biến True & ALO & AMO & Exactly-One\\\midrule
0 & Không thỏa & Thỏa & Không thỏa\\
1 & Thỏa & Thỏa & Thỏa\\
Từ 2 trở lên & Thỏa & Không thỏa & Không thỏa\\
\bottomrule\end{tabular}
\end{table}

ALO trực tiếp có một clause dài $k$, không thêm biến phụ. Những phương pháp trong báo cáo thay đổi phần AMO; khi cần Exactly-One, tất cả cộng cùng một ALO trực tiếp. Đây là cách tách ALO/AMO trong slide bài học \citep{amoSlides}. Một số tài liệu giới thiệu ALO thông qua commander hoặc mã bit, nhưng các lựa chọn đó không phải cài đặt đang được đánh giá.

\subsection{Biến phụ và tính đúng của một encoding}
Gọi $E(X,Y)$ là CNF dùng biến chính $X$ và biến phụ $Y$. Để mã hóa AMO đúng trên đầu vào, cần
\begin{equation}
\AMO(X)\ \Longleftrightarrow\ \exists Y\,E(X,Y).
\label{eq:existential}
\end{equation}
Chiều từ phải sang trái bảo đảm CNF không cho phép hai đầu vào cùng đúng. Chiều từ trái sang phải bảo đảm mọi đầu vào hợp lệ đều có cách gán biến phụ thỏa CNF. Không yêu cầu các biến phụ có cách gán duy nhất. Khi tìm hoặc đếm bố trí hậu, cần chiếu nghiệm về biến ô cờ để không đếm lặp do các cách gán phụ khác nhau.

Kích thước encoding gồm số biến phụ, số clause và số literal. Chúng tác động đến chi phí sinh/nạp và cấu trúc suy diễn, nhưng không tự quyết định thời gian tìm kiếm. Một encoding nhỏ hơn có thể tạo đường tìm kiếm khó hơn đối với một backend cụ thể.

\subsection{SAT solving và các phương pháp CP/ILP}
Glucose3 là backend SAT được dùng trong đợt thực nghiệm. Các bộ giải SAT hiện đại kết hợp gán quyết định, suy diễn đơn vị, phân tích xung đột và học clause. Khi một literal trong clause được gán khiến mọi literal khác sai, literal còn lại phải đúng. Encoding ảnh hưởng các phép suy diễn này, không chỉ ảnh hưởng số lượng clause đầu vào.

Constraint Programming (CP) biểu diễn biến có miền và các ràng buộc mức cao như AllDifferent. Trong N-Queens, ba AllDifferent trên $q_i$, $q_i+i$ và $q_i-i$ mô tả gọn điều kiện không tấn công. CPLEX CP trong báo cáo là IBM CP Optimizer qua \texttt{docplex.cp}.

CP-SAT là bộ giải ràng buộc nguyên của OR-Tools. Nó có kỹ thuật dựa trên SAT nhưng đầu vào trong báo cáo là mô hình nguyên với AllDifferent, nên được đánh giá riêng với SAT CNF tự mã hóa. Không gọi đây là bộ giải CP thuần hoặc đồng nhất ba ràng buộc mức mô hình với ba clause nội bộ.

Integer Linear Programming (ILP) dùng biến nguyên và các ràng buộc tuyến tính. Bài này dùng biến nhị phân cho từng ô; CPLEX MIP và Gurobi giải cùng hệ ràng buộc khả thi. Mặc dù tên backend là MIP, mô hình N-Queens ở đây không có biến liên tục. Các phương pháp đều chính xác theo nghĩa lời giải được kiểm tra và kết quả UNSAT phải là chứng minh không tồn tại nghiệm; timeout không phải chứng minh vô nghiệm.
